% !TEX TS-program = lualatex+se
% !TEX encoding = UTF-8

\documentclass[12pt, a4paper]{extarticle}
%\usepackage[letterpaper,showframe]{geometry}
\usepackage{fullpage} % redukcja marginesów
\usepackage{geometry}
\geometry{
	left=10mm,
	right=10mm,
	top=12mm,
	bottom=12mm,
	marginparwidth=0mm,
	marginparsep=0mm
	}
	
% ------------- FONTY ------------
\usepackage[latin,polish]{babel}
\usepackage{fontspec}
%\usepackage{times}
%\usepackage[math]{anttor}
%\usepackage{tgtermes}
%\usepackage{mathptmx}
\usepackage{ebgaramond}
\setmainfont{EB Garamond}
%\usepackage{libertine}
%\setmainfont{Libertinus Serif}
\defaultfontfeatures{Ligatures=TeX}	
% --- GREGORIO ---	
\usepackage{gregoriotex}
\gresetlinecolor{gregoriocolor}
\gresetheadercapture{annotation}{greannotation}{string}% odczytuje adnotacje nad inicjałem jako ciąg, co pozwala na jego renderowanie
\gresetheadercapture{commentary}{grecommentary}{string}% odczytuje komentarz jako ciąg, co pozwala na jego renderowanie
\grechangedim{commentaryraise}{0.3cm}{scalable} % podnosi komentarz wystarczająco ponad górną linię
\grechangestyle{annotation}{\color{black}\small}% format adnotacji
\grechangestyle{commentary}{\color{gregoriocolor}\small}% format komentarza
\grechangestyle{abovelinestext}{\color{gregoriocolor}\small}% format napisu nad liniami
\grechangestyle{initial}{\fontsize{36}{36}\selectfont} % zmniejsza wielkość inicjału do bardziej proporcjonalnego rozmiaru (=36pt)
\grechangedim{intersyllablespacenotes}{0.26cm}{scalable}% Minimalna odległość między nutami różnych sylab.

% umieszczenie zawartości składu ich unikalnego argumentu po lewej i prawej stronie każdej linii
\let\grelocalleftbox\localleftbox
\let\grelocalrightbox\localrightbox

%\newcommand{\separator}{\greseparator{3}{40}}
\newfontfamily\vectis{../fonts/VectisW01-Bold.ttf}
%\def\greinitialformat#1{\fontsize{20}{20}\vectis #1}
% --- format pierwszej linii 
% \setfirstannotation{\fontsize{10}{10}\textcolor{red}{Ant.}}%
% --- format drugiej linii
% \setsecondannotation{\fontsize{10}{10}\sc{\textcolor{red}8.}}%
% --- dodawanie spacji
% \tolerance=9999 %(or less)
% \pretolerance=500  %default is -1  This will spread your chant out over the bars
% --- wymuś hypheny między sylabami
% \grechangedim{\gremaximumspacewithoutdash}{0pt}
% --- redukcja spacji między sylabami i glyphami: \hspace{-.2ex}
% --- lub wstam sylaby lub glyphy między to: \mbox{...}
% --- lub użyj ( z), aby przesunąć neum bliżej początku pięciolinii. Im więcej spacji w nawiasach, tym bardziej zawęża spację.
% ------------------

%  Konfiguracja kolumn
\usepackage{paracol}
\setcolumnwidth{12.5cm,5.5cm}
\setlength{\columnsep}{0.5cm}
\footnotelayout{m}
% \setlength{\columnseprule}{0.4pt}
% \colseprulecolor{red}
\usepackage{tcolorbox}
\newtcolorbox{nutybox}{size=minimal, colback=white, colframe=red, arc=0mm, width=\linewidth, halign=left, valign=center}
\newtcolorbox{tlumbox}{size=minimal, colback=white, colframe=red, arc=0mm, width=\linewidth, halign=left, valign=top, fontupper=\linespread{.8}\selectfont}

\usepackage{Tabbing}
\usepackage{longtable}
\usepackage[savepos]{zref}

\pagestyle{plain}	% usuwanie nagłówków/stopek z\,dokumentu
\usepackage[indent=0pt]{parskip} % ustawianie paragrafów
\usepackage{setspace} % odstępy między liniami
%\setstretch{1.5} % półtorej linii odległości między wierszami

\usepackage{graphicx} % umożliwia import obrazów
\graphicspath{{img/}} % katalog obrazów
\usepackage{pgfornament}% ornamenty graficzne

\usepackage{titlesec} % ustawianie nagłówków

% --- komendy własne ---
\newcommand{\red}[1]{\textcolor{red}{#1}}
\newcommand{\black}[1]{\textcolor{black}{#1}}

\newcommand{\koniec}{
	\begin{center}
%		\vspace{5mm}
%		\pgfornament[width=0.6cm]{#1}
		\newpage
	\end{center}}

\begin{document}
\titlespacing*{\chapter}{0pt}{0pt}{5pt}	
\titleformat{\chapter}[hang]{\filcenter\LARGE\scshape\bfseries}{}{0pt}{} % format nagłówków
\titlespacing*{\section}{0pt}{0pt}{0pt}	
\titleformat{\section}[hang]{\filcenter\large\scshape}{}{0pt}{} % format nagłówków

\setlength{\parindent}{0pt}
\gresetbracerendering{font}

\pagenumbering{gobble} % nie jest wyprowadzany żaden numer strony, chociaż numer jest nadal resetowany. Odniesienia do tej strony również są puste.
\large
\begin{spacing}{0.8}
\begin{center}
	\section*{Caelitum Joseph decus}
	\small\textit{Liber antiphonarius, Solesmes, 1960, p. 690}
\end{center}	
\end{spacing}
\begin{paracol}{2}
	\ensurevspace{4cm}
	\begin{nutybox}
	\greannotation{Hymnus}\greannotation{1.}\gabcsnippet{(c4) CAé(h)li(g')tum(h) Jo(gf)seph(dc) de(g')cus,(g) at(f')que(g) no(h.)strae(h.) (;) Cer(h)ta(j') spes(k) vi(ji)tae,(h.) (,) co(j')lu(g)mén(f')que(g) mun(h.)di,(h.) (:) Quas(hih) ti(g')bi(h) lae(ef)ti(d.) (,) cá(g')ni(g)mus,(f) be(gh)ní(hgg)gnus(f.) (,) Sús(f!gwh_g)ci(e')pe(f) lau(d.)des.(d.) (::)}
	\end{nutybox}
	\switchcolumn
	\begin{tlumbox}
	\vspace{3mm} 
	\red{1}.~Niebian chlubo, Józefie, i~pewna nadziejo życia naszego, oraz podporo świata, przymij życzliwie pieśni, jakie Ci śpiewamy radośnie.
	\end{tlumbox}
  	\ensurevspace{4cm}
	\switchcolumn*
	\begin{nutybox}
	\gresetinitiallines{0}\gabcsnippet{(c4)  <c>2.</c> Te(h) Sa(g')tor(h) re(gf)rum(dc) stá(g')tu(g)it(f') pu(g)dí(h.)cae(h.) (;) Vír(h)gi(j')nis(k) spon(ji)sum,(h.) (,) vo(j')lu(g)ít(f')que(g) Ver(h.)bi(h.) (:) Te(hih) pa(g')trem(h) di(ef)ci,(d.) (,) de(g')dit(g) et(f) mi(gh)ní(hgg)strum(f.) (,) Es(f!gwh_g)se(e') sa(f)lú(d.)tis.(d.) (::)}
	\end{nutybox}
	\switchcolumn
	\begin{tlumbox}
	\vspace{3mm} 
	\red{2}.~Ciebie Stwórca (wszech) rzeczy postanowił oblubieńcem czystej Dziewicy i~chciał, (by) Cię zwano Ojcem Słowa, dozwolił Ci także być wspólnikiem zbawienia.
	\end{tlumbox}
	\ensurevspace{4cm}
	\switchcolumn*
	\begin{nutybox}
	\gresetinitiallines{0}\gabcsnippet{(c4)  <c>3.</c> Tu(h) Red(g')em(h)ptó(gf)rem(dc) stá(g')bu(g)lo(f') ja(g)cén(h.)tem,(h.) (;) Quem(h) cho(j')rus(k) Va(ji)tum(h.)  (,) cé(j')ci(g)nit(f') fu(g)tú(h.)rum,(h.) (:) A(hih)spi(g')cis(h) gau(ef)dens(d.) (,) hu(g')mi(g)lís(f)que(gh) na(hgg)tum(f.) (,) Nu(f!gwh_g)men(e') ad(f)ó(d.)ras.(d.) (::)}
	\end{nutybox}
	\switchcolumn
	\begin{tlumbox}
	\vspace{3mm} 
	\red{3}.~Ty Odkupiciela, leżącego w~stajence, którego przyjście zapowiadał chór Proroków, oglądasz z~radością i~pokornie oddajesz pokłon Bogu narodzonemu.
	\end{tlumbox}
	\ensurevspace{4cm}
	\switchcolumn*
	\begin{nutybox}
	\gresetinitiallines{0}\gabcsnippet{(c4)  <c>4.</c> Rex(h) De(g')us(h) re(gf)gum,(dc) Do(g')mi(g)ná(f')tor(g) or(h.)bis,(h.) (;) Cu(h)jus(j') ad(k) nu(ji)tum(h.) (,) tre(j')mit(g) in(f')fe(g)ró(h.)rum(h.) (:) Tur(hih)ba,(g') c{ui}(h) pro(ef)nus(d.) (,) fa(g')mu(g)lá(f)tur(gh) ae(hgg)ther,(f.) (,) Se(f!gwh_g) ti(e')bi(f) sub(d.)dit.(d.) (::)}
	\end{nutybox}
	\ensurevspace{5cm}
	\switchcolumn
	\begin{tlumbox}
	\vspace{3mm}
	\red{4}.~Bóg, Król królów, Pan świata, na którego skinienie drży zgraja piekielna, któremu w~pokorze służy niebo, poddaje się Tobie.
	\end{tlumbox}
	\ensurevspace{4cm}
	\switchcolumn*
	\begin{nutybox}
	\gresetinitiallines{0}\gabcsnippet{(c4)  <c>5.</c>  Laus(h) sit(g') ex(h)cél(gf)sae(dc) Trí(g')a(g)di(f') per(g)én(h.)nis,(h.) (;) Quae(h) ti(j')bi(k) prae(ji)bens(h.) (,) sú(j')pe(g)ros(f') ho(g)nó(h.)res,(h.) (:) Det(hih) tu(g')is(h) no(ef)bis(d.) (,) mé(g')ri(g)tis(f) be(gh)á(hgg)tae(f.) (,) Gáu(f!gwh_g)di(e')a(f) vi(d.)tae.(d.) (::) A(ded)men.(cd..) (::)}
	\end{nutybox}
	\switchcolumn
	\begin{tlumbox}
	\vspace{3mm}
	\red{5}.~Chwała ustawiczna niechaj będzie Trójcy najwyższej, która dając Ci zaszczyty niebieskie, niechaj dla twoich zasług udzieli nam radości żywota błogosławionego. Amen.\footnotemark
	\end{tlumbox}	
	\footnotetext{tłum. Ks. Bronisław Gładysz, "Hymny brewiarza rzymskiego oraz patronału polskiego. Przekład i objaśnienia", Poznań 1933} 
\end{paracol}	
\end{document}